\chapter{Análisis del problema}\label{cap:analisis}

\section{Introducción}
En este capítulo se formaliza el problema desde la perspectiva de ingeniería de
datos: qué información se necesita, qué funcionalidades debe cumplir el sistema,
y qué restricciones garantizan que los resultados sean técnicamente defendibles.

El objetivo no es construir un producto comercial, sino una solución académica
completa, reproducible y trazable, capaz de convertir datos masivos de viajes en
evidencia analítica útil.

\section{Preguntas de análisis que guían el proyecto}
El sistema se diseñó para responder, al menos, las siguientes preguntas:
\begin{itemize}
	\item ¿Cómo se distribuye la demanda de viajes por franja horaria y día de semana?
	\item ¿Qué zonas concentran mayor volumen de recogidas y destinos?
	\item ¿Qué diferencias económicas aparecen por zona (tarifa, propina, ingreso)?
	\item ¿Existe un desequilibrio territorial entre recogidas y destinos por distrito?
	\item ¿Cómo difiere el comportamiento de viajes estándar frente a viajes de aeropuerto?
\end{itemize}

Estas preguntas se reflejan directamente en los notebooks de Zeppelin, que actúan
como pieza de síntesis del trabajo.

\section{Requisitos de información}
Los datos base proceden del repositorio público de NYC TLC
(\textit{Yellow Taxi Trip Records}) \cite{nyctlc,nyctlcparquet}. La fuente principal
se consume en formato Parquet por su adecuación a cargas analíticas
\cite{parquetdocs}.

Campos mínimos utilizados:
\begin{itemize}
	\item \texttt{tpep\_pickup\_datetime}, \texttt{tpep\_dropoff\_datetime};
	\item \texttt{passenger\_count}, \texttt{trip\_distance};
	\item \texttt{PULocationID}, \texttt{DOLocationID};
	\item \texttt{payment\_type}, \texttt{fare\_amount}, \texttt{tip\_amount},
	\texttt{total\_amount};
	\item \texttt{VendorID}, \texttt{RatecodeID}, \texttt{store\_and\_fwd\_flag}.
\end{itemize}

Datos geográficos auxiliares:
\begin{itemize}
	\item Shapefile de zonas TLC, convertido a GeoJSON para enriquecer y visualizar
	por \texttt{LocationID}.
\end{itemize}

\section{Requisitos funcionales}
\textbf{RF-01. Carga de datos Parquet.} El sistema debe cargar el fichero
\texttt{yellow\_tripdata\_2025-01.parquet} en Spark y validar su esquema.

\textbf{RF-02. Procesamiento distribuido.} El sistema debe cargar y transformar los
datos con Apache Spark en modo distribuido local.

\textbf{RF-03. Limpieza trazable.} El sistema debe aplicar reglas explícitas sobre
nulos, duración, distancia, importes, pasajeros y coherencia temporal.

\textbf{RF-04. Variables derivadas.} El sistema debe generar variables para análisis
temporal (fecha, hora, día, franja horaria) y operativo (duración, velocidad).

\textbf{RF-05. Traducción de códigos.} El sistema debe convertir codificación numérica
(vendor, pago, tarifa, flag) a etiquetas legibles de negocio.

\textbf{RF-06. Integración geográfica.} El sistema debe enriquecer el dataset por
zonas de origen y destino usando el GeoJSON de zonas TLC.

\textbf{RF-07. Consultas analíticas.} El sistema debe ejecutar consultas por hora,
franja, tipo de viaje, pasajeros, distrito y balance territorial.

\textbf{RF-08. Exportación de indicadores por zona.} El sistema debe generar CSVs de
indicadores (viajes, tarifa, propina, distancia, duración, ingresos, hora pico/valle)
por zona de origen y de destino.

\textbf{RF-09. Visualización geográfica.} El sistema debe generar mapas coropléticos
interactivos con Folium, con inspección por tooltip.

\textbf{RF-10. Trazabilidad de ejecución.} El sistema debe ofrecer salidas verificables
en los notebooks y archivos de indicadores para su revisión posterior.

\section{Requisitos no funcionales}
\textbf{RNF-01. Reproducibilidad.} La ejecución debe replicarse entre equipos mediante
Docker Compose.

\textbf{RNF-02. Claridad técnica.} Notebooks y documentación deben ser
comprensibles para revisión académica.

\textbf{RNF-03. Trazabilidad.} Cada fase del flujo debe dejar evidencias de salida
(recuento de limpieza, indicadores exportados y visualizaciones generadas).

\textbf{RNF-04. Escalabilidad conceptual.} La solución local debe permitir justificar
una migración futura a clúster real.

\textbf{RNF-05. Honestidad de resultados.} No se deben reportar métricas inventadas;
la discusión debe apoyarse en ejecuciones observables.

\textbf{RNF-06. Mantenibilidad.} El repositorio debe separar datos, notebooks y
documentación para facilitar evolución.

\section{Restricciones y supuestos}
\begin{itemize}
	\item El entorno de ejecución es local; por tanto, memoria y CPU son limitadas.
	\item El análisis se centra en un mes de referencia (enero de 2025) para
	mantener tiempos de procesamiento razonables.
	\item La calidad del dato depende de la fuente pública; el sistema no corrige
	errores semánticos externos no detectables por reglas objetivas.
	\item El GeoJSON de zonas se genera mediante un script auxiliar
	(\texttt{convert.py}) a partir del shapefile oficial de TLC.
\end{itemize}

\section{Criterios de aceptación}
Se consideran cumplidos los requisitos cuando:
\begin{itemize}
	\item el notebook principal ejecuta sus bloques de carga, limpieza y análisis
	sin errores;
	\item se obtienen los 14 ficheros CSV de indicadores por zona (7 de origen y
	7 de destino);
	\item el notebook de mapas genera visualizaciones coropléticas navegables;
	\item se documenta el flujo completo con evidencias de ejecución.
\end{itemize}

\section{Conclusiones}
El análisis confirma que el problema es abordable con un flujo de analítica
distribuida en Spark siempre que se mantengan tres principios: limpieza trazable,
reproducibilidad operativa y transparencia en la interpretación de resultados.
Con estos requisitos, los notebooks de Zeppelin se convierten en evidencia
integrada de implementación y análisis.